% No \movedhead here, unlike the other moved sections. body.tex prints this
% heading in both builds, because sec:mediator and sec:lengthorlabel are
% referenced from shared prose and so have to resolve to a Section in both.
% A \movedhead as well set the title twice in the reading build, as an empty
% subsection followed by the real one.
\label{app:mediator}
The argument for setting the reasoning budget through the API rather than
measuring how much reasoning a response contains is made in
Section~\ref{sec:deliberation} on identification grounds: a measure read off
the response is produced by the same forward pass as the answer, so it cannot
separate reasoning that helped from reasoning that a hard item provoked. This
data can test that argument instead of asserting it, because it holds both
variables. The effort parameter is exogenous, set before the model sees the
item. Output tokens are endogenous, chosen by the model for each response.
They are not the self-reported analytical-reasoning indicator used by
\citet{bini2026behavioral}, but they are the same class of variable: a
within-cell, model-chosen measure of how much reasoning a response contains,
with no experimental control over it. None of this subsection is
pre-registered; it was written after the data were in hand.

The parameter does set reasoning volume, so the two variables are not
measuring different constructs. Output tokens rise from \texttt{low} to
\texttt{high} by a factor of 2.33 in Haiku, 2.44 in Opus and 2.91 in Sonnet,
and the distributions barely overlap: 0.1\% of Haiku's and Opus's
\texttt{low} effort responses exceed the median \texttt{high} effort
response, 0.0\% of Sonnet's, and in the other direction 0.0\%, 0.0\% and
3.5\%. If an observational measure of reasoning volume recovers the causal
effect of reasoning volume anywhere, it should do so here.

It does not. \texttt{analysis/mediator.py} runs the paper's own arithmetic
three ways on the same 9{,}720 responses, with the same families, items and
statistics, changing only which responses enter each pair mean.

\paragraph{Exogenous.} Comparing effort levels, the mean absolute family
effect falls from \texttt{low} to \texttt{high} in all three models on both
measures, by $-0.120$, $-0.069$ and $-0.050$ on the categorical measure for
Opus, Sonnet and Haiku, and by $-2.628$, $-3.667$ and $-3.700$ on the
continuous one. Six of six.

\paragraph{Endogenous, between items.} Within a single model and effort cell
the effort label is constant, but generation length still varies widely from
item to item. If deliberation itself is what reduces bias, then holding the
label fixed, pairs on which the model happened to generate more should show
less bias. They show more. The Spearman correlation between a pair's mean
output tokens and its absolute effect, pooled across all nine cells with the
cell means removed from both variables so that only within-cell variation
contributes, is $+0.190$ ($p = 0.0001$) on the categorical measure and
$+0.154$ ($p = 0.0001$) on the continuous measure, over 972 pairs. The
relationship is positive in 17 of the 18 cell by measure combinations and
reaches $p < 0.05$ in eight of them. Taken as a median split instead, so that
the statistic matches the one used for the exogenous comparison, the mean
absolute effect is lower in the heavier half in one cell of nine on both
measures, and pooled over effort and bootstrapped over pairs the contrast is
positive in all three models on both measures, with Haiku on the categorical
measure excluding zero in the wrong direction at $+0.066$ $[+0.014, +0.143]$.
Between cells, bias falls as effort rises. Within cells, bias rises with
length. The two slopes have opposite signs.

\paragraph{Endogenous, within item.} Two readings of that reversal were
available. The first is that the exogenous budget and the model's own
allocation of tokens are distinct interventions with distinct effects. The
second is reverse causation: an item that provokes bias also provokes a
longer answer, so length is an index of item difficulty rather than of
deliberation. A third design separates them. Inside one model by effort cell,
split the five trials of each single item at that item's own median token
count. The stimulus, the family, the company and the effort setting are then
identical across the split, and only the model's sample to sample choice of
how much to write differs. The reversal disappears and nothing takes its
place: the mean absolute effect falls in four of nine cells, which is what
chance gives, and every pooled interval contains zero, the largest being
$-0.585$ $[-2.896, +1.999]$ for Haiku on the continuous measure. The positive
association is therefore a property of the item, not of the response, and the
difficulty reading is the one the data support.

An observational study of these responses would thus have reported that
reasoning makes the bias worse, or that it does nothing, depending on whether
it conditioned on the item. The experiment reports that it reduces the bias,
in every model and on both measures.

Leverage rules out the obvious objection that the observational designs
simply had less variation to work with. Both separate their halves by a
similar amount, 49.5 and 53.6 output tokens in Opus, 59.6 and 61.2 in Sonnet,
132.0 and 155.1 in Haiku, which is between 0.19 and 0.39 of the gap the
parameter opens. The two observational designs therefore have the same
leverage as each other and less than the experiment. Less leverage is a fair
explanation of the within-item null, which should be read as consistent with
low power rather than as evidence against the experiment. It is not an
explanation of a sign reversal. Where the between-item design returns the
wrong sign it is returning confounding, not noise.

Because the two observational designs differ only in whether the item is held
fixed, their disagreement locates that confounding in the item, and two item
level properties move with the token count. Harder items draw more words and
are more bias-prone, which is the reading above. The second is visible in the
data and is less often named: the families themselves differ in how much the
model writes about them, so splitting on tokens re-weights the cues being
averaged. Mental accounting contributes 75.5\% of its responses to the heavy
half and price path 68.2\%, while action bias contributes 28.0\% and framing
30.2\%. A measure of reasoning volume taken from the response is therefore
also a measure of which cue was present.

What no observational design here can remove is the loop that motivates the
study. A response's recommendation and the length of its justification are
produced together, so conditioning on length conditions on the answer.
Setting the budget before the item is seen is the only one of the three
designs that avoids this, and it is the only one that recovers an effect. A
design which measures deliberation by observing the model, whether through
generated length or through the model's own report of its reasoning, is
measuring a quantity whose association with bias runs the other way from the
one an experimenter can set. That is not a refinement of the identification
problem in \citet{bini2026behavioral}; it is a demonstration of it. It is not
a demonstration that their finding is wrong: their mediator, items and models
all differ from these. It shows that the class of design cannot be relied on
to recover the sign, which is a claim about method and not about their
result.
